\documentclass[12pt]{article}

% Layout.
\usepackage[top=1.2in, bottom=0.9in, left=1in, right=1in, headheight=1in, headsep=6pt]{geometry}

% Fonts.
\usepackage{mathptmx}
\usepackage[scaled=1.0]{helvet}
\renewcommand{\emph}[1]{\textsf{\textbf{#1}}}

% Misc packages.
\usepackage{amsmath,amssymb,latexsym}
\usepackage{graphicx,hyperref}
\usepackage{array}
\usepackage{xcolor}
\usepackage{multicol}
\usepackage{tabularx,colortbl}
\usepackage{enumitem}

\hypersetup{
    colorlinks=true,
    linkcolor=blue,
    filecolor=magenta,      
    urlcolor=blue,
    pdftitle={Syllabus for MATH F252 Spring 2026, in-person sections},
    pdfpagemode=FullScreen,
    }

\def\mailto#1{\href{mailto:#1}{#1}}

% Paragraph spacing
\parindent 0pt
\parskip 6pt plus 1pt
\def\tableindent{\hskip 0.5 in}
\def\ts{\hskip 1.5 em}

\usepackage{fancyhdr}
\pagestyle{fancy} 
\lhead{\large\sf\textbf{MATH F252 Calc 2}}
\chead{\large\sf\textbf{Syllabus (section 001)}}
\rhead{\large\sf\textbf{Spring 2026}}

\newcommand{\localhead}[1]{\par\smallskip\textbf{#1} \smallskip\nobreak\\}%
\def\heading#1{\localhead{\large\emph{#1}}}
\def\subheading#1{\localhead{\emph{#1}}}

\newenvironment{clist}%
{\bgroup\parskip 0pt\begin{list}{$\bullet$}{\partopsep 4pt\topsep 0pt\itemsep -2pt}}%
{\end{list}\egroup}%

\begin{document}

\strut\par\vskip-12pt
\heading{Essential Information}

\vskip -12pt
\renewcommand{\arraystretch}{1.25}
\begin{tabular}{lllcl}
& instructor: &Glen Woodworth\\
\\
& section: & 001 (CRN 30347) \\
& class times: & Asynchronous\\
& classrooms: & None\\
\\
& instructor offices: & Zoom only\\
& office hours: & M 3:30-4:30 Zoom \cr
& &T 3:45-4:45 Zoom \cr 
& &Th 10:00-11:00 Zoom\cr \\
& emails: & \href{mailto:gwoodworth@alaska.edu}{\texttt{gwoodworth@alaska.edu}} \\
\end{tabular}
 
 \renewcommand{\arraystretch}{1.5}
\begin{tabular}{ll}
{\emph{Course webpages}} & public: \href{https://uaf-math.github.io/calc2/}{\texttt{uaf-math.github.io/calc2}}\\
 & canvas (section 001): \href{https://canvas.alaska.edu/courses/30029}{\texttt{canvas.alaska.edu/courses/30029}} \\

\emph{Prerequisite} & MATH F251X Calculus I; or placement. \\
{\emph{Required text}} &\textit{OpenStax Calculus Volume 2} by G. Strang \& E. Herman,\cr
&\href{https://openstax.org/details/books/calculus-volume-2}{\texttt{openstax.org/details/books/calculus-volume-2}} \\
&optional paperback print copy:\quad \texttt{ISBN-13 978-1-50669-807-6} 
\end{tabular}

\heading{Description and Course Goals}
Calculus is part of the language in all the science and engineering disciplines.  It is \textsl{useful}, and that is why it is part of the UAF core curriculum.  The two principal tools from Calculus I (MATH 251), differentiation and integration, are central here too.  However, this course extends your understanding of integration, and we will also study sequences and series, including Taylor series.

At the completion of the course, students will:
\begin{clist}
\item possess mature skills in computing integrals in one variable,
\item be able to apply integrals to common problems (areas, volumes, work, \dots),
\item understand the convergence of sequences and series,
\item understand and be able to apply Taylor series and polynomials, and
\item be able to use parametric and polar equations for curves.
\end{clist}
In addition, students will have the mathematical foundation for success in Calculus III and other courses requiring knowledge of sophisticated integration techniques in one variable, sequences and series, including many junior/senior-level science and engineering courses.

\clearpage\newpage
\phantom{x}
\vspace{-4mm}

\heading{Learning Calculus in an asynchronous online class}
This is an asynchronous distance course. Instead of attending in-person classes, students will be introduced to new material by reading the text and by viewing pre-recorded videos posted on Canvas. Each week a student  will be reading 2-3 sections from the text with 3-5 videos consisting of 1-2 hours of content. Students should budget 3-4 hours to read and watch the
videos so that they can take notes and reread or rewatch parts that did not make sense.

\heading{Schedule and Online Materials}
The course website contains a \href{https://uaf-math.github.io/calc2/assets/general/S26/asynchronous/schedule_v1.pdf}{Schedule} listing the textbook sections to be covered each class, the dates each Homework is due, plus the dates for Quizzes and Exams. You should consult this schedule frequently.  We will announce any Schedule adjustments in class.

Most course materials (Syllabus, old Quizzes, Exams with solutions, study materials, etc.) will be posted on the \href{https://uaf-math.github.io/calc2/}{public course webpage}.  A few course materials (grades, Homework solutions, announcements, link to Gradescope) will be available on the Canvas sites; see the first page of this Syllabus.

\heading{Office Hours and Communication}
Canvas will be used to make announcements for the whole class.  (\textsl{Please set up your Canvas accounts to receive these announcements in a location that is regularly checked!})  See the office hours on the first page of this syllabus.  You are also welcome to make an appointments by email to the instructor email addresses on the first page.\\


\heading{Evaluation and Grades}
Grades are determined as follows.  (Each component of the grade is discussed below.)
 
\begin{multicols}{2}
\begin{tabular}{|c|c|}
\hline
Weekly Check-in & 5\% \\
\hline
Homework & 15\% \\
\hline
Quizzes & 15\% \\
\hline
Midterm Exam 1 & 20\% \\
\hline
Midterm Exam 2 & 20\%  \\
\hline
Final Exam & 25\% \\
\hline
total & 100\% \, \\
\hline
\end{tabular}

%\vskip 6pt

\begin{tabular}{llll}
A  & 93--100\%& C  & 70--76\%  \\
A- & 90--92\% & &  \\
B+ & 87--89\% & D+ & 67--69\%  \\
B  & 83--86\% & D  & 63--66\%  \\
B- & 80--82\% & D- & 60--62\%  \\
C+ & 77--79\% & F  & $\le$ 59\%
\end{tabular}
\end{multicols}

The grade ranges at right are a guarantee and a lower bound. We reserve the right to increase your grade above these ranges based on the actual difficulty of the work, average class performance, and/or improvement over the semester including exemplary performance on the final exam. \\

\clearpage\newpage
\phantom{x}
\vspace{-4mm}

\heading{Homework}
Homework assignments consist of a selection of problems from the textbook.  All problems for the entire semester can be found at the public webpage \href{https://uaf-math.github.io/calc2/homework.html}{here}. 

The Homework exists because all of us \textbf{learn by doing}.

Please write your Homework on paper, or electronically on a tablet, and turn it in as a PDF via Gradescope, accessed via the Canvas pages (see first page of this Syllabus).  

Help with scanning homework can be found on the \href{https://uaf-math.github.io/calc2/techHelp.html}{Tech Help} webpage.  Assignments are due, mostly on Mondays and Wednesdays, by 11:59pm, thus in advance of the Thursday Quiz; see the online \href{https://uaf-math.github.io/calc2/assets/general/S26/inperson/schedule.pdf}{Schedule} for due dates.  

Complete worked solutions to all Homework problems are \textbf{provided in advance on the Canvas site for your section}.  Thus, Homework will be graded based on {effort} and {completion}.  However, \textbf{the work you turn in is expected to be your own}.  

To earn full points on the homework, your submission must:
\vskip -20pt\strut
\begin{clist}
\item be unequivocally your own work
\item include work, not bald answers
\item be legible, organized, and in order
\item be on time
\end{clist}

\heading{Weekly Check-ins}
At the beginning of every week, there is a Check-in assignment. The goal of these is three-fold: (1) to ensure a student is aware of all the deadlines in the week ahead, (2) to provide an avenue for regular communication with the instructor and (3) to provide a mechanism for a student to trouble-shoot any challenges. These will be graded on completion. 

\heading{Quizzes}
On most weeks, there will be a quiz.  The quiz consists of problems like the homework problems due that week. The purpose of the quiz is to help you consolidate your learning from the week, to identify areas of confusion, and then provide an opportunity to fix any misconceptions.  

You will download the quiz from Gradescope, take the quiz on your own without any aids. Once you attempted all the problems on your own, you will download the solutions to the quiz and \textcolor{red}{correct your quiz in red pen}. You will upload your corrected quiz. 

\textbf{Your quiz work will earn full points provided that your initial work is your own and that you fully correct your paper, thus all students should earn 100\% on the quizzes.} Note that the consequence of this grading scheme is to establish that there is no benefit to your grade by using external sources when initially working the quiz problems and substantial disincentive for violating this rule. \textbf{Specifically, to use external sources when you initially complete the quiz constitutes a violation of academic integrity and the Student Code of Conduct. }

For all quizzes, you will print the blank quiz or download the quiz to a tablet. Complete the quiz on your own using no aids. 

\clearpage\newpage
\phantom{x}
\vspace{-4mm}

\heading{Midterm and Final Exams}
There are two Midterm Exams and a Final Exam this semester, to be held on the dates in the schedule on the course website. The Final Exam is cumulative. All Midterms and the Final Exam will be closed book, closed notes, and no calculator. All Midterms and the Final Exam will be proctored.\\
%\vskip -20pt\strut

\begin{clist}
\item \emph{Midterm 1, Thursday February 19 or Friday February 20}
\item \emph{Midterm 2, Thursday April 9 or Friday April 10}
\item  \emph{Final Exam, Wednesday April 29 or Thursday April 30 or Friday May 1}
\end{clist} 

\heading{A Typical Week (with suggested timing)}
\vspace*{-.2in}
\begin{enumerate}
\item Complete the Check-in. (Monday)
\item Finish homework, check answers, ask questions. Submit last week's homework. (Monday) 
\item Read the text and/or watch the video introducing the first new section. Start homework. (Monday)
\item Finish last week's homework problems, check answers, ask questions.  (Tuesday)
\item Read/Watch intro to the second new section. Submit homework. (Wednesday)
\item Start homework. (Thursday)
\item Take the quiz, correct your answers. Submit your corrected quiz.  (Thursday/Friday)
\item Start the next Section. (Friday)
\end{enumerate}

\heading{Proctored Assessments}
There are \emph{three} proctored assessments for this class:
Midterm 1, Midterm 2, and the Final Exam. 

\heading{Scheduling your proctored assessments through eCampus}
For the proctored assessments, you will set up the proctoring arrangement through eCampus by following these instructions.
\begin{itemize}
\item If you live in the Fairbanks area, you can schedule your proctored assessments by going
the eCampus Exam Services site \href{https://ecampus.uaf.edu/exam-services/}{https://ecampus.uaf.edu/exam-services/} and clicking
the yellow box labeled “Schedule a Testing Appointment” near the middle of the page.
For the exam group, select “UAF eCampus Course Proctored Exam”. Select “Paper Test”
and choose the appropriate exam from our course, Math F252X.
\item If you live outside the Fairbanks area, you can schedule your proctored assessments by
going the eCampus Exam Services site \href{https://ecampus.uaf.edu/exam-services/}{https://ecampus.uaf.edu/exam-services/} and
scrolling down to the menu of “Popular Exams” where you will find a yellow bar labeled
“UAF eCampus Courses”. Read this section until you see two options in blue lettering: “I
am in Alaska (but not in Fairbanks)” and ”I am outside of Alaska.” Select the option that
applies to you. When completing your form, select “UAF eCampus Course Proctored
Exam”, “Paper Test” and choose the appropriate exam from our course, Math F252X.
\end{itemize}

\heading{Getting Help}
\vspace*{-.3in}

Calculus 2 is a challenging course!  (You will earn your sense of accomplishment when you succeed at learning the material.)  Please expect to experience the discomfort that is an unavoidable part of learning something new, including periods of confusion and struggle. You should expect to need to get extra help from your instructor, tutors in the Math and Stat Lab, and each other. 

Please come and talk to us if you are unable to clarify points of confusion. We want to know what your long-term struggles are and help you move past them.

Places to Get Help

\vskip -20pt\strut
\begin{clist}
    \item The Math and Stat Lab is now located in the
      Student Success Center on the 6th floor of the Rasmusen Building and
      offers drop-in tutoring.

	See \href{http://www.uaf.edu/dms/mathlab/}{\texttt{www.uaf.edu/dms/mathlab/}} for schedules and availability.
	\item Free
one-on-one (or small group) tutoring is also available in 
the SSC. You must schedule an
appointment; see \href{http://www.uaf.edu/dms/mathlab/}{\texttt{www.uaf.edu/dms/mathlab/}}.
	\item Student Support Services (\href{https://uaf.edu/sss/}{\texttt{uaf.edu/sss/}}) offers free tutoring in many subjects to students who qualify for their program.
	\item ASUAF (\href{https://uaf.edu/asuaf/}{\texttt{uaf.edu/asuaf/}}) offers private tutoring for a small fee, based on student income.
\end{clist}


\clearpage\newpage
\phantom{x}
\vspace{-4mm}

\heading{Rules and Policies}
\vskip -20pt

\subheading{AI usage}
Feel free to use a calculator or online tools on Homework.  It is also reasonable to explore new AI tools like ChatGPT, but merely doing cut-and-paste without understanding will have no benefit to your learning. 

Moreover, during proctored and on-paper Quizzes and Exams, electronic tools of any type are not allowed.  Since on-paper Quizzes and Exams represent 85\% of your grade, copying without understanding is not a good long-term strategy.

\subheading{Incomplete Grade} 
Incomplete (I) will only be given in
  DMS courses in cases where
  the student has completed the majority (normally all but the last
  three weeks) of a course with a grade of C or better, but for
  personal reasons beyond his/her control has been unable to complete
  the course during the regular term. Negligence or indifference are
  not acceptable reasons for the granting of an incomplete. 

\subheading{Late Withdrawals} 
A withdrawal after the deadline from a DMS course will
  normally be granted only in cases where the student is performing
  satisfactorily (i.e., C or better) in a course, but has exceptional
  reasons, beyond his/her control, for being unable to complete the
  course. These exceptional reasons should be detailed in writing to
  the instructor, department head and dean.

\subheading{No Early Final Examinations}
Final examinations for DMS
  courses shall not be held earlier than the date and time published
  in the official term schedule. Normally, a student will not be
  allowed to take a final exam early. Exceptions can be made by
  individual instructors, but should only be allowed in exceptional
  circumstances and in a manner which doesn't endanger the security of
  the exam.

\subheading{Academic Dishonesty}
Academic dishonesty, including cheating and plagiarism, will not
be tolerated.  It is a violation of the Student Code of Conduct
and will be punished according to UAF procedures.


\subheading{Student protections and services statement}
Every qualified student is welcome in my classroom.  As needed, I am happy to work with you, Disability Services, Veterans' Services, Rural Student Services, etc.~to find reasonable accommodations.  Students at this University are protected against sexual harassment and discrimination (Title IX), and minors have additional protections.  As required, if I notice or am informed of certain types of misconduct, then I am required to report it to the appropriate authorities.  For more information on your rights as a student and the resources available to you to resolve problems, please go the following site: \href{https://www.uaf.edu/orca/index.php}{\texttt{https://www.uaf.edu/orca/index.php}}.

%\clearpage\newpage

\strut

\vspace{-12pt}

\subheading{General education statement}
This course is listed as a General Education Math Course.  As such this course is expected to \textsl{contribute to meeting the following} four general learning outcomes:

\begin{enumerate}
\item Build knowledge of human institutions, sociocultural processes, and the physical and natural works through the study of mathematics.  Competence will be demonstrated for the foundational information in each subject area, its context and significance, and the methods used in advancing each.

\item Develop intellectual and practical skills across the curriculum, including inquiry and analysis, critical and creative thinking, problem solving, written and oral communication, information literacy, technological competence, and collaborative learning. Proficiency will be demonstrated across the curriculum through critical analysis of proffered information, well-reasoned solutions to problems or inferences drawn from evidence, effective written and oral communication, and satisfactory outcomes of group projects.

\item Acquire tools for effective civic engagement in local through global contexts, including ethical reasoning, intercultural competence, and knowledge of Alaska and Alaska issues.  Facility will be demonstrated through analyses of issues including dimensions of ethics, human and cultural diversity, conflicts and interdependencies, globalization, and sustainability.   

\item Integrate and apply learning, including synthesis and advanced accomplishment across general and specialized studies, adapting them to new settings, questions and responsibilities, and forming a foundation for lifelong learning. Preparation will be demonstrated though production of a a creative or scholarly product that requires broad knowledge, appropriate technical proficiency, information collection, synthesis, interpretation, presentation and reflection.
\end{enumerate}

%\newpage

\begin{center}
\textbf{\large{Official UAF Syllabus Addendum}}
\end{center}


\noindent{\bf Student protections statement:} The university respects and upholds the principles of due process and a fair and equitable process as specified in the Board of Regents' Policy 09.02 Student Rights and Responsibilities. For more information regarding the rights and responsibilities of students, refer to the Office of Rights, Compliance and Accountability website. You are encouraged to read the Board of Regents' policy carefully to fully understand your responsibilities to our community.

We strive to create a safe and respectful environment for all members of our community. If you have questions about expectations of you as a student or believe your rights are being violated, we encourage you to reach out to the  Office of Rights, Compliance and Accountability for help. UAF reserves the right to suspend, expel or take other necessary and appropriate action in cases where a student is unable or unwilling to uphold community standards and campus safety.

For more information on your rights as a student and the resources available to you to resolve problems, please go to the following site:\\ \url{https://catalog.uaf.edu/academics-regulations/students-rights-responsibilities/}

\noindent{\bf Disability services statement:} I will work with the Office of Disability Services to provide reasonable accommodation to students with disabilities.

\noindent{\bf ASUAF advocacy statement:} The Associated Students of the University of Alaska Fairbanks, the student government of UAF, offers advocacy services to students who feel they are facing issues with staff, faculty, and/or other students specifically if these issues are hindering the ability of the student to succeed in their academics or go about their lives at the university. Students who wish to utilize these services can contact the Student Advocacy Director by visiting the ASUAF office or emailing asuaf.office@alaska.edu. 



\noindent{\bf Student Academic Support:}
\begin{itemize}
\setlength\itemsep{0em}
        \item Communication Center (907-474-7007, \mailto{uaf-commcenter@alaska.edu}, Student Success Center, 6th Floor Room 677 Rasmuson Library)
        \item Writing Center (907-474-5314, \mailto{uaf-writing-center@alaska.edu}, Student Success Center, 6th Floor Room 677 Rasmuson Library)
\item UAF Math Services (907-474-7332, \mailto{uaf-traccloud@alaska.edu})


\begin{itemize}
\item Drop-in tutoring, Student Success Center, 6th Floor Room 672 Rasmuson Library

\item 1:1 tutoring (by appointment only), 6th Floor Room 677 Rasmuson Library

\item Online tutoring (by appointment only) available

https://www.uaf.edu/dms/mathlab/, available at the Student Success Center
\end{itemize}

\item Developmental Math Lab, Gruening 406
\item The Debbie Moses Learning Center at CTC (907-455-2860, 604 Barnette St, Room 120,\\ \url{https://www.ctc.uaf.edu/student-services/student-success-center/})
\item For more information and resources, please see the Academic Advising Resource List (\url{https://www.uaf.edu/advising/students/index.php})
\end{itemize}

\noindent{\bf Student Resources:}
\begin{itemize}
\setlength\itemsep{0em}
\item Disability Services (907-474-5655, \mailto{uaf-disability-services@alaska.edu}, 110 Eielson Building)
\item Student Health \& Counseling [free counseling sessions available] (907-474-7043, \url{https://www.uaf.edu/chc/appointments.php}, Whitaker Building, Room 206, Health, Safety \& Security Bldg --- same building as Fire and Police)
\item Office of Rights, Compliance and Accountability (907-474-7300, \mailto{uaf-orca@alaska.edu}, 3rd Floor, Constitution Hall)
\item Associated Students of the University of Alaska Fairbanks (ASUAF) or ASUAF Student Government (907-474-7355, \mailto{asuaf.office@alaska.edu}{asuaf.office@alaska.edu}, Wood Center 119)
\end{itemize}

\noindent{\bf Nondiscrimination statement:}
Nondiscrimination statement: The University of Alaska is an equal opportunity/equal access employer, educational institution and provider. The University of Alaska does not discriminate on the basis of race, religion, color, national origin, citizenship, age, sex, physical or mental disability, status as a protected veteran, marital status, changes in marital status, pregnancy, childbirth or related medical conditions, parenthood, sexual orientation, gender identity, political affiliation or belief, genetic information, or other legally protected status. The University's commitment to nondiscrimination, including against sex discrimination, applies to students, employees, and applicants for admission and employment. Contact information, applicable laws, and complaint procedures are included on UA's statement of nondiscrimination available at \url{www.alaska.edu/nondiscrimination}.

\begin{tabular}{l}
UAF Office of Rights, Compliance and Accountability\\
1692 Tok Lane\\
3rd floor, Constitution Hall, Fairbanks, AK 99775\\
907-474-7300\\
\url{uaf-orca@alaska.edu}
\end{tabular}

\hfill  \scriptsize [syllabus version 1.1: \today] \normalsize

\end{document}
